\section{Arithmetic and the theory of groups}
\label{sec:annex-p02-group-theory-undecidability}

\begin{annexitem}{annex:p02.finite-ring-group-condition}{thm:p02-finite-ring-group-condition}{Ring condition}
\begin{proof}
Assume \(\Theta_0(a,b)\). Its first clause says that every commutator is
central. In such a group
\[
 [xx',y]=[x,y][x',y],\qquad [x,yy']=[x,y][x,y'].
\]
For example, use \(xy=yx[x,y]\) to move \(y\) past \(x'\) and then
\(x\); centrality lets the two commutators move past all remaining
factors. Comparing with \((xx')y=y(xx')[xx',y]\) and cancelling proves
the first identity. Moving \(x\) past \(y\) and then \(y'\) proves
the second. Taking one of the products to be the identity also gives
the corresponding inverse identities.

Put \(X=C_G(a)\), \(Y=C_G(b)\), and \(Z=Z(G)\).
Centralizers and the centre are subgroups, by multiplying and inverting
their defining commutation equalities. The centre is abelian. The maps
\[
 X\longrightarrow Z,\quad x\longmapsto[x,b],\qquad
 Y\longrightarrow Z,\quad y\longmapsto[a,y]
\]
are homomorphisms by the displayed identities and are surjective by
\(\beta(a,b)\). Each has kernel \(Z\): its kernel is the intersection
\(C_G(a)\cap C_G(b)\), which is the centre by the second clause of
\(\Theta_0\).

For central \(z,w\), surjectivity supplies witnesses \(x\in X\),
\(y\in Y\) with \([x,b]=z\) and \([a,y]=w\). Their commutator is
central and supplies a value of \(M_{a,b}(z,w,v)\). If \(x'\) is
another such witness for \(z\), then \(x^{-1}x'\) lies in the first
kernel and is central. Likewise two witnesses for \(w\) differ by
a central element. Multiplying either argument of a commutator by a
central element leaves that commutator unchanged. Thus \(v=[x,y]\)
does not depend on the witnesses. The graph is total and single-valued
on the centre, with central values.

Write this operation as \(z*w\). Representatives \(x_1,x_2\) of
\(z_1,z_2\) have product \(x_1x_2\) representing \(z_1z_2\).
The commutator identities give
\[
 (z_1z_2)*w=(z_1*w)(z_2*w),\qquad
 z*(w_1w_2)=(z*w_1)(z*w_2).
\]
Hence it distributes over the additive group, whose operation here is
written multiplicatively. The central element \(c=[a,b]\) is a
two-sided multiplicative identity: choose \(x=a\) when the first
argument is \(c\), and \(y=b\) when the second is \(c\).
The two extra clauses in \(\Theta_{\mathrm{ring}}\) make \(*\)
commutative and associative. For associativity, totality supplies the
three products in the antecedent of its displayed clause; the consequent
identifies their result with the other parenthesization. Therefore the
centre, with these operations, is a commutative ring with identity.
Its additive inverse is the group inverse. The domain is nonempty,
equality is literal, and all operation graphs are total and functional;
these also verify the structural requirements for the interpretation
\(I\), not just the ring identities.

Conversely take \(\mathrm H(R)\) with its standard pair
\(e_1=(1,0,0)\), \(e_2=(0,1,0)\).
By \cref{thm:p03-w03-p03-074-heisenberg-basic}, every commutator is
central. For \(g=(u,v,w)\), commuting with \(e_1\) means \(v=0\),
and commuting with \(e_2\) means \(u=0\). Their common centralizer
is exactly \(\zeta(R)=Z(\mathrm H(R))\). For a central element
\(\zeta(r)\), the witnesses \((r,0,0)\) and \((0,r,0)\) give the
two required commutators, so \(\beta(e_1,e_2)\) holds. Their determinant is one.
Finally \cref{thm:p03-w03-p03-074-ring-recovery}, with determinant one,
identifies the interpreted ring with \(R\). Its multiplication is
therefore associative and commutative, proving the remaining clauses.
All calculations include the zero ring.
\end{proof}
\end{annexitem}

\begin{annexitem}{annex:p02.finite-group-arithmetic}{thm:p02-finite-group-arithmetic}{The arithmetic interpretation}
\begin{proof}
All component formulas are finite, and \(A_{\mathsf Q}\) has seven
members. Substituting their fixed translations into \(\vartheta\), and
then expanding the fixed interpretation formulas, gives a finite group
formula \(\Theta(a,b)\). Thus \(V_G\) is a finite first-order axiom set
in a finite signature. No arithmetic operation is added as a primitive
group symbol.

Let \((G,c_x,c_y)\models V_G\). By
\cref{thm:p02-finite-ring-group-condition}, \(I(G,c_x,c_y)\) is a
commutative ring \(R\). It satisfies \(\vartheta\), by the truth
property of the interpretation. In particular \(D(R)\) contains
zero and one and is closed under addition and multiplication. On this
domain take arithmetic zero to be ring zero, successor to be addition
of ring one, and the other operations to be the restrictions of the
ring operations. They are total on the domain.

For any assignment in \(D(R)\), induction on arithmetic terms shows
that their values are the values of their translated ring terms.
Closure supplies the successor, addition, and multiplication cases.
Induction on formulas then proves that arithmetic truth in this
structure is equivalent to truth of \(J(\psi)\) in \(R\): atomic
equality is unchanged, Boolean connectives commute with translation,
and the quantified clauses restrict exactly to \(D(R)\).
Since \(R\models J(\alpha)\) for every axiom \(\alpha\) of
\(\mathsf Q\), this arithmetic structure is a model of \(\mathsf Q\).

Composing with \(I\), the arithmetic domain in \(G\) is the definable
subset of its centre satisfying \(D^I\). Equality remains literal.
Arithmetic zero is the group identity, successor multiplies a central
element by \([c_x,c_y]\), addition is group multiplication on the
domain, and multiplication is the restriction of
\(\mathrm{Mul}(c_x,c_y;-,-,-)\). The closure sentence proves that the
outputs stay in the arithmetic domain. The graph translation for these
data has the truth property proved by induction above. Consequently it
agrees in every \(V_G\)-model with \((J(\psi))^I\), and completeness
gives the displayed provable equivalence. This also combines the nested
centre and four-square quantifier restrictions into the single domain
guard of \(K\), without identifying two differently written formulas.
Each structural interpretation condition and each translated arithmetic
axiom holds in every \(V_G\)-model; strong first-order completeness
makes it a \(V_G\)-theorem. Thus this is an interpretation in the
syntactic sense of \cref{def:p05-relative-interpretation}, as well as
a construction of interpreted structures. The constants \(c_x,c_y\)
belong to its target signature and are not free parameters in its data.

The translation is effective. The fixed arithmetic-to-ring translation
is primitive recursive on the specified codes by
\cref{thm:p02-w02c-p02-003-ring-coding}; the ring-to-group translation
replaces symbols by finitely specified domain and operation-graph
formulas. The same finite graph-recursion procedure applied to the
displayed composite data produces \(K\). Its recursion on a finite
formula, including fresh-variable renaming, is effective.

For consistency, use the single model
\((\mathrm H(\mathbb Z),e_1,e_2)\).
Its interpreted ring is \(\mathbb Z\) by
\cref{thm:p02-finite-ring-group-condition}. In that ring every sum of
four squares is nonnegative, and every nonnegative integer is such a
sum by \cref{thm:p02-w02c-p02-003-lagrange}. Thus \(D(\mathbb Z)\)
is exactly the set of natural numbers. It contains zero and one, is
closed under the two operations, and satisfies the seven Robinson
axioms with the stated successor. Consequently
\(\mathbb Z\models\vartheta\), and the truth equivalence for \(I\)
gives \((\mathrm H(\mathbb Z),e_1,e_2)\models V_G\). Its interpreted
arithmetic is the standard one.
Soundness proves consistency of \(V_G\).
\end{proof}
\end{annexitem}

\begin{annexitem}{annex:p02.group-theory-undecidable}{thm:p02-group-theory-undecidable}{Decision reduction}
\begin{proof}
By \cref{thm:p05-inc-theoremhood-chain}, every consistent c.e.
extension of \(\mathsf Q\) has undecidable theoremhood; the displayed
certificate interfaces there have established instances for each such
extension. Thus \(\mathsf Q\) is essentially undecidable in the sense
of \cref{def:p05-essential-undecidability}.
The finite consistent theory \(V_G\) effectively interprets it by
\cref{thm:p02-finite-group-arithmetic}. The proved interpretation-transfer
theorem \cref{thm:p05-interpretation-transfer} makes \(V_G\)
essentially undecidable, and in particular undecidable.

Suppose there were a decider for sentences true in all groups. Given
an arbitrary sentence \(\sigma\) of the expanded group signature of
\(V_G\), choose variables \(a,b\) not occurring in it, and replace the
two extra constants by those variables, obtaining \(\sigma(a,b)\).
Compute the pure group sentence
\[
 \forall a\,\forall b\,
       \bigl(\Theta(a,b)\to\sigma(a,b)\bigr).
\]
It is true in every group exactly when \(\sigma\) is true in every
\(V_G\)-model: an assignment satisfying \(\Theta\) is precisely an
expansion by the two constants satisfying \(V_G\), and assignments
not satisfying \(\Theta\) make the implication true. By strong
first-order completeness the latter condition is exactly
\(V_G\derives\sigma\). The assumed decider would therefore decide
theoremhood in \(V_G\), a contradiction. This is the finite-axiom
step needed to pass from the interpreted subclass to the theory of all
groups; undecidability of one particular group's complete theory alone
would not supply it.

Finally, the group axioms form a finite effective theory. Enumerating
finite first-order proofs from them enumerates their sentence
consequences. Soundness and completeness identify those consequences
with the sentences true in all groups. Hence that set is c.e. as well
as undecidable.
\end{proof}
\end{annexitem}
